% =====================================================================
% Report content to be pasted INSIDE \begin{document} ... \end{document}
%
% Required packages in your Overleaf preamble:
%   \usepackage{graphicx}
%   \usepackage{amsmath, amssymb}
%   \usepackage{float}        % for [H] figure placement
%   \usepackage{booktabs}     % nice tables (\toprule, \midrule, \bottomrule)
%   \usepackage{listings}     % code snippets
%   \usepackage{xcolor}
%   \usepackage{hyperref}     % cross-references
%   \lstset{basicstyle=\ttfamily\footnotesize, breaklines=true,
%           keywordstyle=\color{blue}, commentstyle=\color{gray!80!black},
%           numbers=left, numberstyle=\tiny, stepnumber=1, frame=single,
%           language=Matlab}
%
% Place the folder 'image_project4' (with sub-folders ex1, ex2, ex3 and
% the file sim_general.png) next to your main.tex
%
% NOTE: three figures of Exercise 2 have a dot in the file name
% (sigma_r_-0.1.png, ...). They are included with the {<name>}.png brace
% trick so that graphicx parses the extension correctly.
% =====================================================================

\title{Anti-lock Braking System Design \\ \large Project 4 -- Driver Assistance System Design (01USHLO)}
\author{Dennis Bettinsoli \\ Politecnico di Torino}
\date{\today}
\maketitle

%======================================================================
\section{Introduction}
%======================================================================
The braking performance of a road vehicle is limited by the friction
that the tyres can exchange with the road, which in turn depends
strongly on the \emph{longitudinal slip} of each wheel. When the brake
torque is too high the wheel decelerates faster than the vehicle, the
slip grows, the wheel locks and two undesirable effects occur
simultaneously: the available longitudinal force drops below its peak
value and the lateral (cornering) force collapses, so the vehicle loses
both deceleration capability and steerability. The purpose of an
Anti-lock Braking System (ABS) is precisely to keep the slip close to
the value that maximises the road friction, thus minimising the
stopping distance while preserving directional control.

This project is organised in three exercises, all built on the same
longitudinal driveline/vehicle model implemented in Simulink:
\begin{itemize}
    \item \textbf{Exercise 1 -- Driveline model (braking).} Open-loop
          braking manoeuvres from $80$\,km/h: the influence of the
          brake command $h_b$ and of the front/rear brake ratio $K_b$
          on the stopping distance, the sensitivity to vehicle mass and
          drag coefficient in a $\pm15\%$ range, and the comparison
          between a geared stop and a neutral-gear stop ($h_c=0$).
    \item \textbf{Exercise 2 -- Slip-regulation ABS.} Design and tuning
          of a closed-loop wheel-slip controller
          ($\dot u = K_\mathrm{gain}\,\mathrm{sign}(\sigma_r-\sigma)$,
          $T_\mathrm{brake}=\mathrm{sat}(u,0,T_b^{(f,r)})$): choice of
          the reference slip $\sigma_r$, gain tuning for minimum
          stopping distance, behaviour on a wet road and a discussion of
          its practical feasibility.
    \item \textbf{Exercise 3 -- Conjugate Boundary Method.} A
          slip-free ABS based on a Stateflow state machine that
          modulates the brake pressure using only the wheel
          acceleration, tested on wet and snow surfaces.
\end{itemize}

The nominal vehicle parameters (front-wheel/rear-wheel drive selectable
through the flag \texttt{Traction}) are summarised in
Table~\ref{tab:params}; the initial speed is $V_0=80/3.6\approx
22.2$\,m/s in every test.

\begin{table}[H]
    \centering
    \begin{tabular}{lcl}
        \toprule
        \textbf{Symbol} & \textbf{Value} & \textbf{Description} \\
        \midrule
        $m$        & $1020$\,kg     & Nominal vehicle mass \\
        $C_x$      & $0.33$         & Nominal drag coefficient \\
        $a$        & $0.923$\,m     & Front axle -- c.o.g. distance \\
        $b$        & $1.376$\,m     & Rear axle -- c.o.g. distance \\
        $R_e$      & $0.279$\,m     & Effective rolling radius \\
        $T_{b,\mathrm{tot}}$ & $3000$\,Nm & Total available brake torque \\
        $V_0$      & $80$\,km/h     & Initial speed \\
        $T_s$      & $1$\,ms        & Solver / discrete sample time \\
        \bottomrule
    \end{tabular}
    \caption{Main vehicle parameters used throughout Project 4.}
    \label{tab:params}
\end{table}

%======================================================================
\section{Theoretical Background}
\label{sec:theory}
%======================================================================

\paragraph{Longitudinal slip.}
For a braked wheel the longitudinal slip is defined as
\begin{equation}
    \sigma = \frac{R_e\,\omega - V}{V},
    \label{eq:slip}
\end{equation}
where $\omega$ is the wheel angular speed and $V$ the vehicle forward
speed. During braking $R_e\omega < V$, hence $\sigma<0$; $\sigma=0$ is
free rolling and $\sigma=-1$ is a fully locked wheel ($\omega=0$). This
sign convention explains why all the reference slips and controller
gains used in the code are negative.

\paragraph{Friction curve.}
The road force is modelled through a $\mu$--$\sigma$ characteristic
(stored as a 1-D look-up table). The friction coefficient grows almost
linearly with $|\sigma|$, reaches a peak $\mu_\mathrm{max}$ at an
optimal slip $\sigma_p$ (typically $|\sigma_p|\approx0.1\!-\!0.2$ on
dry asphalt, lower on wet/snow) and then \emph{decreases} towards the
sliding value as $|\sigma|\to1$:
\begin{equation}
    F_x = \mu(\sigma)\,F_z .
    \label{eq:fx}
\end{equation}
The ascending branch ($|\sigma|<|\sigma_p|$) is \emph{stable}: a small
increase of slip raises the force that tends to spin the wheel back up.
The descending branch ($|\sigma|>|\sigma_p|$) is \emph{unstable}: more
slip means less force, the wheel decelerates further and runs away to
lock. Any ABS must therefore keep the operating point on, or just
below, the peak.

\paragraph{Wheel and vehicle dynamics.}
Each wheel obeys
\begin{equation}
    J_w\,\dot\omega = R_e\,F_x - T_\mathrm{brake},
    \label{eq:wheel}
\end{equation}
while the longitudinal motion of the vehicle is governed by
\begin{equation}
    m\,\dot V = -\!\!\sum_{i} F_{x,i} - \tfrac12\rho\,C_x A\,V^2 - F_\mathrm{roll},
    \label{eq:long}
\end{equation}
i.e. the sum of the four tyre forces plus the aerodynamic drag and the
rolling resistance. Equation~\eqref{eq:long} shows that the mass scales
the inertia (longer stop when heavier) whereas the drag coefficient
only enters the aerodynamic term, which at $\le80$\,km/h is small
compared with the braking forces.

\paragraph{Brake distribution.}
The total brake torque is split front/rear through the ratio $K_b$:
\begin{equation}
    T_{bf}^{\max} = \frac{T_{b,\mathrm{tot}}}{1+1/K_b}, \qquad
    T_{br}^{\max} = \frac{T_{bf}^{\max}}{K_b},
    \label{eq:Kb}
\end{equation}
and the instantaneous command is $T_\mathrm{brake}=h_b\,T_{b}^{\max}$
with $h_b\in[0,1]$. During braking the load transfers forward, so the
rear axle is lightly loaded and locks first; a larger $K_b$ (more front
bias) is therefore safer, while a small $K_b$ promotes the
directionally-unstable rear lock.

%======================================================================
\section{Exercise 1 -- Driveline Model (Braking)}
%======================================================================

\subsection*{Simulink implementation}
The longitudinal model (Figure~\ref{fig:sim_general}) chains the
\texttt{Engine + clutch}, \texttt{Gearbox} and \texttt{Differential}
blocks into the \texttt{Wheel dynamics} subsystem, which integrates
equation~\eqref{eq:wheel} for each wheel, evaluates the slip
\eqref{eq:slip} and the tyre force \eqref{eq:fx} through the
$\mu$--$\sigma$ table, and returns the longitudinal force $F_x$. The
net force is divided by the mass and integrated twice
(equation~\eqref{eq:long}) to obtain speed and travelled distance; a
\texttt{STOP} block ends the run when the speed falls below
$\approx1$\,m/s. The brake torque enters the wheel block scaled by
$h_b$, and a switch (highlighted later in Figure~\ref{fig:wheel_model})
selects between the fixed open-loop torque used here and the controller
output used in Exercise 2.

\begin{figure}[H]
    \centering
    \includegraphics[width=\textwidth]{image_project4/sim_general.png}
    \caption{Top-level longitudinal driveline/vehicle model used in
    Exercise 1: engine--clutch--gearbox--differential feeding the wheel
    dynamics, with double integration to speed and distance.}
    \label{fig:sim_general}
\end{figure}

%----------------------------------------------------------------------
\subsection{Effect of $h_b$ and $K_b$ on the stopping distance}
%----------------------------------------------------------------------
Two extreme calibrations are compared. The first
(Figure~\ref{fig:hbkb1}) uses a gentle command $h_b=0.25$ with a strong
front bias $K_b=2.33$ ($\approx70/30$); the second
(Figure~\ref{fig:hbkb2}) uses a hard command $h_b=0.90$ with a more
rearward split $K_b=1.7$.

\begin{figure}[H]
    \centering
    \includegraphics[width=0.9\textwidth]{image_project4/ex1/h_b_K_b_stoppingDistance_1.png}
    \caption{Gentle braking: $h_b=0.25$, $K_b=2.33$. Speed/distance
    (top) and tyre slip (bottom).}
    \label{fig:hbkb1}
\end{figure}

\begin{figure}[H]
    \centering
    \includegraphics[width=0.9\textwidth]{image_project4/ex1/h_b_K_b_stoppingDistance_2.png}
    \caption{Hard braking: $h_b=0.90$, $K_b=1.7$. The four wheels lock
    ($\sigma=-1$) within $0.1$\,s.}
    \label{fig:hbkb2}
\end{figure}

With $h_b=0.25$ the deceleration is mild: the slip of all wheels stays
close to zero (only a brief rear-right excursion to $\approx-0.45$ near
the end, when the speed is low and the same torque produces more slip),
the vehicle does not even come to a complete stop within the
$4$\,s horizon and the travelled distance reaches
$\approx47$\,m. With $h_b=0.90$ the brake torque immediately exceeds
the grip of every tyre: all four slips saturate at $\sigma=-1$ in less
than $0.1$\,s and stay locked until the vehicle stops at
$\approx35$\,m in $\approx3.05$\,s.

The two cases summarise the role of the two parameters:
\begin{itemize}
    \item \textbf{$h_b$} sets the magnitude of the brake torque, hence
          the deceleration. Increasing $h_b$ shortens the stopping
          distance, \emph{but only until the wheels lock}: beyond the
          lock threshold the slip saturates at $-1$ and any further
          increase of $h_b$ is useless. The shorter distance of the
          hard stop ($35$ vs $47$\,m) is obtained at the price of fully
          locked wheels -- i.e. zero steering authority and operation on
          the descending (low-$\mu$) branch of the friction curve,
          which on a real road would also lengthen the stop and cause
          loss of control.
    \item \textbf{$K_b$} sets the front/rear distribution
          (equation~\eqref{eq:Kb}). A larger $K_b$ moves torque to the
          (more loaded) front axle and protects the rear from locking;
          the smaller $K_b=1.7$ of the second case sends comparatively
          more torque to the unloaded rear axle and contributes to the
          simultaneous lock.
\end{itemize}
The ideal stop would keep $\sigma$ at the friction peak on every wheel
-- which is exactly what the ABS of Exercises 2 and 3 automates.

%----------------------------------------------------------------------
\subsection{Sensitivity to vehicle mass and drag coefficient ($\pm15\%$)}
%----------------------------------------------------------------------
Keeping a fixed calibration ($h_b=0.45$, $K_b=2.33$) the mass and the
drag coefficient are perturbed together by $\pm15\%$ around their
nominal values ($m_0=1020$\,kg, $C_{x0}=0.33$). The heavy/high-drag
configuration ($+15\%$: $m=1173$\,kg, $C_x=0.3795$) is shown in
Figure~\ref{fig:mass1}, the light/low-drag one ($-15\%$, parameters at
$0.85$ of nominal) in Figure~\ref{fig:mass2}.

\begin{figure}[H]
    \centering
    \includegraphics[width=0.9\textwidth]{image_project4/ex1/vehicleMass_Cx_1.png}
    \caption{Sensitivity, $+15\%$ case ($m=1173$\,kg, $C_x=0.3795$):
    longer stop, rear axle locks at $\approx0.6$\,s.}
    \label{fig:mass1}
\end{figure}

\begin{figure}[H]
    \centering
    \includegraphics[width=0.9\textwidth]{image_project4/ex1/vehicleMass_Cx_2.png}
    \caption{Sensitivity, $-15\%$ case (mass and drag reduced to
    $0.85$ of nominal): shorter stop, rear axle locks earlier
    ($\approx0.45$\,s).}
    \label{fig:mass2}
\end{figure}

In both cases the front wheels keep rolling at a small slip
($\sigma\approx-0.1$) while the lightly-loaded rear axle reaches
$\sigma=-1$. The trend with the parameters is clear and monotonic:
\begin{itemize}
    \item the heavy/high-drag vehicle stops in $\approx33$\,m
          ($\approx2.9$\,s) and the rear locks at $\approx0.6$\,s;
    \item the light/low-drag vehicle stops in $\approx30$\,m
          ($\approx2.6$\,s) and the rear locks earlier, at
          $\approx0.45$\,s.
\end{itemize}
A heavier vehicle therefore needs a longer distance, even though its
higher $C_x$ adds aerodynamic drag: at these speeds the drag term in
equation~\eqref{eq:long} is small, so the \emph{mass dominates} the
sensitivity. The lighter vehicle, having less rear-axle load, reaches
the rear-lock condition sooner. Overall a $\pm15\%$ parameter spread
maps into roughly $\pm5\%$ of stopping distance around the nominal --
a moderate but non-negligible sensitivity that any robust ABS tuning
must tolerate.

%----------------------------------------------------------------------
\subsection{Geared stop versus neutral gear ($h_c=0$)}
%----------------------------------------------------------------------
Finally the same full-pedal stop ($h_b=1$, $K_b=2.33$, nominal vehicle)
is run with the gear engaged ($h_c=1$, engine braking transmitted
through the driveline) and in neutral ($h_c=0$, clutch disengaged, no
engine braking). The two runs are overlaid in
Figure~\ref{fig:neutral}.

\begin{figure}[H]
    \centering
    \includegraphics[width=0.9\textwidth]{image_project4/ex1/neutral_gear.png}
    \caption{Full-pedal stop ($h_b=1$): geared and neutral
    speed/distance curves are practically superimposed
    ($\approx35$\,m), because the wheels lock immediately.}
    \label{fig:neutral}
\end{figure}

The speed and distance curves of the two configurations are practically
indistinguishable: both stop at $\approx35$\,m in $\approx3.05$\,s and
all wheels lock ($\sigma=-1$) within the first tenth of a second. The
reason is physical: once a wheel is locked its angular speed is zero and
its deceleration is set entirely by the tyre--road sliding friction,
not by the driveline. Engine braking can add deceleration \emph{only
while the driven (rear) wheels are still rolling}; under a hard,
lock-up stop that contribution vanishes, so the gear has a negligible
effect on the stopping distance. A measurable geared-vs-neutral
difference would appear only with a gentler $h_b$ that keeps the wheels
rolling, where engine braking on the rear axle would slightly increase
the deceleration of the geared case.

%======================================================================
\section{Exercise 2 -- ABS Controller for Wheel-Slip Regulation}
%======================================================================

\subsection*{Controller implementation}
The open-loop brake torque of Exercise 1 is replaced, on each wheel, by
the closed-loop law requested in the assignment
\begin{equation}
    \dot u = K_\mathrm{gain}\,\mathrm{sign}(\sigma_r-\sigma), \qquad
    T_\mathrm{brake} = \mathrm{sat}\!\left(u,\,0,\,T_b^{(f,r)}\right).
    \label{eq:abs2}
\end{equation}
Inside the wheel subsystem (Figure~\ref{fig:wheel_model}) a switch
selects, for each corner, the controller output instead of the fixed
torque. The per-wheel signal chain (Figure~\ref{fig:wheel_chain})
computes the static slip, passes it through a tyre
\texttt{relaxation\_length} block, reads the friction from the
$\mu$--$\sigma$ table and feeds the slip to the
\texttt{wheel\_braking\_controller}. The controller itself
(Figure~\ref{fig:controller}) is the literal block-diagram translation
of equation~\eqref{eq:abs2}: the error $\sigma_r-\sigma$ enters a
\texttt{sign} block, is multiplied by $K_\mathrm{gain}$, integrated by
$1/s$ and clipped by a saturation to the admissible torque range
$[0,\,T_b^{(f,r)}]$.

\begin{figure}[H]
    \centering
    \includegraphics[width=\textwidth]{image_project4/ex2/simulink_of_theWheels_model.png}
    \caption{Wheel dynamics with the brake-torque switch (highlighted):
    open-loop torque $T_{bf}^{\max}h_b$ or controller output
    $T_\mathrm{brake}$.}
    \label{fig:wheel_model}
\end{figure}

\begin{figure}[H]
    \centering
    \includegraphics[width=0.9\textwidth]{image_project4/ex2/simulink_wheel_brk_controller.png}
    \caption{Per-wheel chain: static slip $\to$ relaxation length $\to$
    $\mu$--$\sigma$ table, with the slip routed to the braking
    controller.}
    \label{fig:wheel_chain}
\end{figure}

\begin{figure}[H]
    \centering
    \includegraphics[width=0.9\textwidth]{image_project4/ex2/wheel_barking_controller.png}
    \caption{Wheel braking controller implementing
    $\dot u=K_\mathrm{gain}\,\mathrm{sign}(\sigma_r-\sigma)$ followed by
    the saturation $\mathrm{sat}(u,0,T_b^{(f,r)})$.}
    \label{fig:controller}
\end{figure}

\paragraph{Sign of the gain.}
With $\sigma,\sigma_r<0$, if the wheel is braking less than the target
($\sigma>\sigma_r$) the error $\sigma_r-\sigma$ is negative, so a
\emph{negative} $K_\mathrm{gain}$ produces $\dot u>0$, increases the
brake torque and drives $\sigma$ towards $\sigma_r$. This is why the
code uses $K_\mathrm{gain}=-600$ (and more negative values when tuning).

%----------------------------------------------------------------------
\subsection{Baseline run and slip analysis ($K_\mathrm{gain}=-600$)}
%----------------------------------------------------------------------
The baseline manoeuvre ($\sigma_r=-0.2$, $K_\mathrm{gain}=-600$,
$h_b=0.35$, $K_b=2.33$) is reported in the $\sigma_r=-0.2$ panel of
Figure~\ref{fig:sigma02}. The integral-of-sign law drives the regulated
(rear) slip towards $-0.2$ and then holds it in a bounded
\emph{limit-cycle} oscillation -- the steady-state signature of the
relay (\texttt{sign}) nonlinearity. No wheel locks; the front wheels
settle at a smaller slip ($\approx-0.05$) because the forward load
transfer lets them carry the brake torque without reaching the target.
The stopping distance is $\approx60$\,m in $\approx4.25$\,s: long,
because with $K_\mathrm{gain}=-600$ the torque integrator ramps slowly
and the average deceleration is modest. The gain is tuned in
Section~5.3.

%----------------------------------------------------------------------
\subsection{Choice of the reference slip $\sigma_r$}
%----------------------------------------------------------------------
The reference slip is the operating point the ABS aims at on the
$\mu$--$\sigma$ curve. The three values $\sigma_r=-0.1,\,-0.2,\,-0.3$
(all with $K_\mathrm{gain}=-600$) are compared in
Figures~\ref{fig:sigma01}--\ref{fig:sigma03}.

\begin{figure}[H]
    \centering
    \includegraphics[width=0.85\textwidth]{{image_project4/ex2/sigma_r_-0.1}.png}
    \caption{$\sigma_r=-0.1$: slip stays shallow and smooth, but the
    grip is under-used (operating point below the friction peak).}
    \label{fig:sigma01}
\end{figure}

\begin{figure}[H]
    \centering
    \includegraphics[width=0.85\textwidth]{{image_project4/ex2/sigma_r_-0.2}.png}
    \caption{$\sigma_r=-0.2$: the slip oscillates in a bounded limit
    cycle around the target, close to the friction peak -- best
    compromise.}
    \label{fig:sigma02}
\end{figure}

\begin{figure}[H]
    \centering
    \includegraphics[width=0.85\textwidth]{{image_project4/ex2/sigma_r_-0.3}.png}
    \caption{$\sigma_r=-0.3$: the target lies on the unstable branch;
    the rear slip repeatedly dives to $-1$ (intermittent lock) and
    recovers.}
    \label{fig:sigma03}
\end{figure}

\begin{itemize}
    \item \textbf{$\sigma_r=-0.1$} -- the target sits on the steep,
          \emph{stable} part of the curve, well below the peak. The slip
          is smooth and free of oscillation, but the friction
          coefficient is not fully exploited, so the deceleration is
          lower and the stop is the longest ($\approx61$\,m). It is the
          safest but least efficient choice.
    \item \textbf{$\sigma_r=-0.2$} -- the target is at/near the peak of
          the friction curve. The slip oscillates in a small bounded
          limit cycle around $-0.2$, the friction coefficient is close
          to $\mu_\mathrm{max}$ and the stop is slightly shorter
          ($\approx60$\,m) at the same gain. This is the best trade-off
          between braking force and stability, which is why $-0.2$ is
          selected.
    \item \textbf{$\sigma_r=-0.3$} -- the target is \emph{beyond} the
          peak, on the unstable descending branch
          (Section~\ref{sec:theory}). As soon as the slip exceeds the
          peak the friction drops, the wheel decelerates further and
          runs away: the rear slip dives to $-1$ for sustained intervals
          (around $t\approx2.4\!-\!3$\,s and again near $t\approx4.2$\,s)
          before the controller manages to recover it. This intermittent
          locking gives no distance benefit ($\approx61$\,m) while
          losing braking force and steerability.
\end{itemize}
The slip trends confirm the rule: the reference must be placed at the
friction peak; a smaller magnitude wastes grip, a larger one
destabilises into locking.

%----------------------------------------------------------------------
\subsection{Gain tuning for minimum stopping distance}
%----------------------------------------------------------------------
Keeping $\sigma_r=-0.2$, the gain magnitude is increased so that the
torque integrator reaches the regulating level sooner.
Figure~\ref{fig:tune} shows the result for $K_\mathrm{gain}=-1450$.

\begin{figure}[H]
    \centering
    \includegraphics[width=0.9\textwidth]{image_project4/ex2/tune_K_gain.png}
    \caption{Tuned controller, $\sigma_r=-0.2$, $K_\mathrm{gain}=-1450$:
    the stopping distance drops to $\approx46$\,m while the slip still
    limit-cycles around the target without locking.}
    \label{fig:tune}
\end{figure}

Raising $|K_\mathrm{gain}|$ from $600$ to $1450$ shortens the stop from
$\approx60$\,m ($4.25$\,s) to $\approx46$\,m ($\approx3.5$\,s): the
faster integrator builds the brake torque more quickly, so the wheels
reach the high-friction slip region earlier. The slip keeps oscillating
in a bounded cycle around $-0.2/-0.3$ but never locks. Increasing the
gain further would enlarge the limit-cycle amplitude and risk
overshooting past the friction peak into locking, so
$K_\mathrm{gain}=-1450$ is a good practical minimum-distance setting on
dry road.

%----------------------------------------------------------------------
\subsection{Wet-road conditions and re-tuning}
%----------------------------------------------------------------------
The $\mu$--$\sigma$ block is switched to the wet characteristic and the
gain is re-tuned. Figure~\ref{fig:drywet} compares the dry calibration
($K_\mathrm{gain}=-1450$) with the wet one ($K_\mathrm{gain}=-700$), at
$h_b=0.5$.

\begin{figure}[H]
    \centering
    \includegraphics[width=\textwidth]{image_project4/ex2/dry_vs_wet_road_conditions.png}
    \caption{Dry (left, $K_\mathrm{gain}=-1450$) versus wet (right,
    $K_\mathrm{gain}=-700$) braking with $\sigma_r=-0.2$: speed/distance
    (top) and slip (bottom).}
    \label{fig:drywet}
\end{figure}

Three observations follow:
\begin{itemize}
    \item \textbf{Distance.} The lower wet friction lengthens the stop
          from $\approx39$\,m ($\approx3.0$\,s) on dry to
          $\approx53$\,m ($\approx3.9$\,s) on wet -- the unavoidable
          consequence of the reduced $\mu_\mathrm{max}$.
    \item \textbf{Re-tuning is necessary.} With the aggressive dry gain
          the dry slip already shows a \emph{growing} limit cycle
          (bottom-left); on the lower-grip wet surface the same gain
          would overshoot the (lower, sharper) friction peak even more
          easily and drive the wheel towards locking. Reducing the gain
          to $-700$ yields a well-damped slip that settles cleanly at
          $\sigma_r=-0.2$ (bottom-right).
    \item \textbf{Surface dependence.} The optimal gain is therefore
          surface-dependent: a single fixed gain tuned on dry is too
          aggressive on wet. This lack of robustness of the
          fixed-gain slip controller motivates the threshold-based logic
          of Exercise 3.
\end{itemize}

%----------------------------------------------------------------------
\subsection{Feasibility on a real vehicle}
%----------------------------------------------------------------------
The controller of equation~\eqref{eq:abs2} regulates the slip
$\sigma=(R_e\omega-V)/V$, which requires the true vehicle ground speed
$V$. While the wheel speed $\omega$ is readily measured by wheel-speed
sensors, $V$ is \emph{not} directly measurable precisely when it
matters: during hard braking the wheels slip, so the wheel speed no
longer equals the vehicle speed. Obtaining $V$ would need additional
sensors (optical ground-speed sensor, GPS, inertial integration) or a
dedicated state estimator, so the slip must be \emph{estimated} rather
than measured. In addition, the \texttt{sign}-based law produces
chattering -- continuous high-frequency brake modulation -- which
stresses the hydraulic valves and degrades comfort.

For these reasons this exact strategy cannot be installed as-is.
Its practical relevance is nonetheless high: it defines the ideal
control objective (hold the slip at the friction peak), it is a valuable
reference/benchmark for simulation and Hardware-in-the-Loop testing, and
it makes explicit why production ABS infers the wheel state from the
wheel \emph{acceleration} instead of the slip -- which is exactly the
approach of Exercise 3.

%======================================================================
\section{Exercise 3 -- ABS, Conjugate Boundary Method}
%======================================================================

\subsection*{Control logic and implementation}
The conjugate-boundary ABS avoids any slip measurement and acts only on
the wheel acceleration $\dot\omega$. It is implemented as the Stateflow
machine of Figure~\ref{fig:stateflow}, embedded in the wiring of
Figure~\ref{fig:ex3_sim}: the requested torque $T_\mathrm{req}=
T_{bf}^{\max}h_b$, the wheel acceleration $\dot\omega$, the previous
output (through a unit delay $1/z$, giving $T_\mathrm{in}$) and a
\texttt{wheel\_lock} flag ($=1$ when $|\omega|<0.1$) feed the
controller, whose output $T_\mathrm{out}$ drives the brake.

\begin{figure}[H]
    \centering
    \includegraphics[width=0.95\textwidth]{image_project4/ex3/stateflow.png}
    \caption{Stateflow machine: IDLE, RELEASE and REAPPLY with the
    corresponding output-torque updates and transition guards.}
    \label{fig:stateflow}
\end{figure}

\begin{figure}[H]
    \centering
    \includegraphics[width=0.85\textwidth]{image_project4/ex3/ex3_sim.png}
    \caption{ABS-controller wiring: requested torque, wheel
    acceleration $\dot\omega$, delayed torque $T_\mathrm{in}$ and the
    wheel-lock flag drive the state machine.}
    \label{fig:ex3_sim}
\end{figure}

The three states implement
\begin{align}
    \text{IDLE:} \quad   & T_\mathrm{out} = T_\mathrm{req}, \\
    \text{RELEASE:} \quad & T_\mathrm{out}(k) = T_\mathrm{out}(k-1) - K_\mathrm{release}\,T_s, \\
    \text{REAPPLY:} \quad & T_\mathrm{out}(k) = T_\mathrm{out}(k-1) + K_\mathrm{reapply}\,T_s,
\end{align}
with the guards
\begin{itemize}
    \item $\;\to$ \textbf{RELEASE} when $-\dot\omega > K_1$ (the wheel
          decelerates too fast) \emph{or} \texttt{wheel\_lock};
    \item RELEASE $\to$ \textbf{REAPPLY} when $\dot\omega > K_3$ (the
          wheel spins back up) \emph{or} \texttt{after}$(K_2,\mathrm{sec})$
          (time-out).
\end{itemize}
The system thus dumps brake pressure as soon as a skid is detected and
rebuilds it -- more slowly -- once the wheel recovers, producing the
characteristic saw-tooth pressure modulation that keeps the slip near
the friction peak \emph{without ever computing it}. Following the
assignment hints, the triggers are started from $K_1\approx K_3$, the
charge/discharge rates are kept constant within each scenario, and
$K_\mathrm{reapply}<K_\mathrm{release}$ to avoid chattering. The two
surfaces use the settings of Table~\ref{tab:ex3par}.

\begin{table}[H]
    \centering
    \begin{tabular}{lccc}
        \toprule
        \textbf{Parameter} & \textbf{Wet} & \textbf{Snow} & \textbf{Unit} \\
        \midrule
        $K_1$ (accel.\ trigger) & $40$   & $40$   & m/s$^2$ \\
        $K_3$ (decel.\ trigger) & $35$   & $35$   & m/s$^2$ \\
        $K_2$ (max time RELEASE) & $0.06$ & $0.06$ & s \\
        $K_\mathrm{release}$    & $4500$ & $5500$ & Nm/s \\
        $K_\mathrm{reapply}$    & $500$  & $200$  & Nm/s \\
        $h_b$ / $K_b$           & $0.6/1.8$ & $0.5/1.8$ & -- \\
        \bottomrule
    \end{tabular}
    \caption{Conjugate-boundary ABS settings for the two surfaces.}
    \label{tab:ex3par}
\end{table}

%----------------------------------------------------------------------
\subsection{Wet-surface results}
%----------------------------------------------------------------------
\begin{figure}[H]
    \centering
    \includegraphics[width=0.9\textwidth]{image_project4/ex3/ex3_wet_slip.png}
    \caption{Wet surface: speed/distance (top) and tyre slip (bottom).
    The slip is held in the useful band after a brief initial spike.}
    \label{fig:wet_slip}
\end{figure}

\begin{figure}[H]
    \centering
    \includegraphics[width=0.9\textwidth]{image_project4/ex3/ex3_wet_torque.png}
    \caption{Wet surface: brake-torque saw-tooth on the four wheels
    (front $\sim$ up to $1100$\,Nm, rear $\sim$ up to $600$\,Nm).}
    \label{fig:wet_torque}
\end{figure}

On the wet road the vehicle stops in $\approx4.6$\,s over
$\approx45$\,m. After a short initial dip to $\sigma\approx-0.5$ -- the
transient before the state machine catches the first skid -- the slip
of every wheel is held in the useful band ($\approx-0.1$ to $-0.3$),
oscillating around the friction peak instead of locking. The brake
torque (Figure~\ref{fig:wet_torque}) shows the expected saw-tooth:
a fast drop in RELEASE followed by a slower ramp in REAPPLY, repeated as
a limit cycle; the more loaded front wheels carry the higher torque
($\sim1100$\,Nm peak) than the rear ($\sim600$\,Nm).

%----------------------------------------------------------------------
\subsection{Snow-surface results}
%----------------------------------------------------------------------
\begin{figure}[H]
    \centering
    \includegraphics[width=0.9\textwidth]{image_project4/ex3/ex3_snow_slip.png}
    \caption{Snow surface: speed/distance (top) and tyre slip (bottom).
    Very long stop, slip still kept bounded by the modulation.}
    \label{fig:snow_slip}
\end{figure}

\begin{figure}[H]
    \centering
    \includegraphics[width=0.9\textwidth]{image_project4/ex3/ex3_snow_torque.png}
    \caption{Snow surface: finer and more frequent torque modulation at
    lower levels (front $\sim$ up to $950$\,Nm, rear $\sim100\!-\!150$\,Nm).}
    \label{fig:snow_torque}
\end{figure}

On snow the very low friction stretches the stop to $\approx13.7$\,s and
$\approx140$\,m. Despite the extreme conditions the controller still
prevents sustained lock: after an initial front excursion to
$\sigma\approx-0.75$ the slip is brought back and modulated around
$-0.1$ to $-0.4$ for the rest of the manoeuvre. The torque
(Figure~\ref{fig:snow_torque}) cycles more frequently and at much lower
levels than on wet; the faster discharge ($K_\mathrm{release}=5500$) and
the gentler recharge ($K_\mathrm{reapply}=200$) are precisely what keeps
the slippery wheel from re-locking and avoids chattering.

%----------------------------------------------------------------------
\subsection{Comments}
%----------------------------------------------------------------------
Using only the wheel acceleration and a lock flag, the conjugate-boundary
ABS keeps the slip in the high-friction region on both low-$\mu$
surfaces -- no slip estimation is required, which is the whole point of
the method and its key practical advantage over the controller of
Exercise 2. The stopping distance scales directly with the surface
friction (snow $\gg$ wet), and the saw-tooth brake-torque trace
reproduces the pressure modulation of real production ABS. The price to
pay is a coarser, oscillatory regulation around the peak rather than the
precise slip tracking of the (non-deployable) Exercise-2 controller.

%======================================================================
\section{Benchmark and Summary}
\label{sec:bench}
%======================================================================
Table~\ref{tab:bench} collects the stopping distances of all the
manoeuvres. The figures are read from the speed/distance histories;
distances of different exercises are not directly comparable because the
brake command $h_b$ differs, but within each block the trends are
consistent.

\begin{table}[H]
    \centering
    \footnotesize
    \begin{tabular}{llccl}
        \toprule
        \textbf{Exercise / case} & \textbf{Settings} &
        \textbf{Dist.\ [m]} & \textbf{Time [s]} & \textbf{Wheel behaviour} \\
        \midrule
        Ex1 -- gentle      & $h_b{=}0.25,\,K_b{=}2.33$ & $\approx47$ & $>4$    & no lock (rolls) \\
        Ex1 -- hard        & $h_b{=}0.90,\,K_b{=}1.7$  & $\approx35$ & $3.05$  & all wheels locked \\
        Ex1 -- $+15\%$     & $m{=}1173,\,C_x{=}0.38$   & $\approx33$ & $2.9$   & rear lock @ $0.6$\,s \\
        Ex1 -- $-15\%$     & $0.85\times$ nominal      & $\approx30$ & $2.6$   & rear lock @ $0.45$\,s \\
        Ex1 -- geared/neutral & $h_b{=}1$              & $\approx35$ & $3.05$  & identical, all locked \\
        \midrule
        Ex2 -- $\sigma_r{=}{-}0.1$ & $K_\mathrm{gain}{=}{-}600$  & $\approx61$ & $4.35$ & shallow, stable \\
        Ex2 -- $\sigma_r{=}{-}0.2$ & $K_\mathrm{gain}{=}{-}600$  & $\approx60$ & $4.25$ & limit cycle @ peak \\
        Ex2 -- $\sigma_r{=}{-}0.3$ & $K_\mathrm{gain}{=}{-}600$  & $\approx61$ & $4.35$ & intermittent lock \\
        Ex2 -- tuned       & $\sigma_r{=}{-}0.2,\,K_\mathrm{gain}{=}{-}1450$ & $\approx46$ & $3.5$ & best dry slip ABS \\
        Ex2 -- dry         & $K_\mathrm{gain}{=}{-}1450,\,h_b{=}0.5$ & $\approx39$ & $3.0$ & limit cycle (growing) \\
        Ex2 -- wet         & $K_\mathrm{gain}{=}{-}700,\,h_b{=}0.5$  & $\approx53$ & $3.9$ & well-damped @ $-0.2$ \\
        \midrule
        Ex3 -- wet         & conjugate boundary        & $\approx45$  & $4.6$  & bounded, no lock \\
        Ex3 -- snow        & conjugate boundary        & $\approx140$ & $13.7$ & bounded, no lock \\
        \bottomrule
    \end{tabular}
    \caption{Stopping-distance benchmark across the three exercises.}
    \label{tab:bench}
\end{table}

%======================================================================
\section{Conclusions}
%======================================================================
The project followed the braking problem from the open-loop driveline
model to two complementary ABS designs.
\begin{itemize}
    \item \textbf{Exercise 1} quantified the open-loop behaviour:
          increasing $h_b$ shortens the stop only until the wheels lock,
          a larger $K_b$ (front bias) protects the rear axle, the
          stopping distance is dominated by the vehicle mass (with a
          $\pm15\%$ swing giving $\approx\pm5\%$ distance) and weakly by
          the drag, and -- under a hard lock-up stop -- the engine
          braking of the engaged gear makes no practical difference with
          respect to neutral.
    \item \textbf{Exercise 2} showed that a slip-regulating controller
          recovers the lost braking by holding $\sigma$ at the friction
          peak. The reference $\sigma_r=-0.2$ is the right target
          (below it the grip is under-used, beyond it the wheel locks),
          increasing $|K_\mathrm{gain}|$ minimises the distance up to
          the onset of locking, and the optimal gain is
          surface-dependent (it had to be reduced from $-1450$ on dry to
          $-700$ on wet). The strategy is, however, not directly
          deployable because the slip is not measurable on a real
          vehicle.
    \item \textbf{Exercise 3} delivered a deployable ABS that uses only
          the wheel acceleration. The IDLE/RELEASE/REAPPLY state machine
          modulates the brake pressure in a saw-tooth limit cycle and
          keeps the slip in the useful band on both wet
          ($\approx45$\,m) and snow ($\approx140$\,m) surfaces, with no
          slip estimation and no sustained lock.
\end{itemize}
The overall message is consistent with real ABS practice: locking the
wheels is never optimal -- the shortest \emph{and} controllable stop is
obtained by keeping the longitudinal slip near the friction peak. A
slip controller defines that ideal target but needs an unmeasurable
quantity, whereas the acceleration-based conjugate-boundary method
approximates the same behaviour with quantities that are actually
available on board, which is why it is the approach adopted in practice.
